\chapter*{Preface}
\addcontentsline{toc}{chapter}{Preface}

This book is a comprehensive companion to Stanford's CS336: Language Modeling from Scratch. It follows the course's ``from scratch'' philosophy---building every component of a modern language model, from tokenization to alignment, with full mathematical rigor and detailed implementation guidance.

\section*{Who This Book Is For}

You are a computer science undergraduate (junior or senior level) who:
\begin{itemize}
  \item Can program fluently in Python
  \item Has taken introductory courses in machine learning and linear algebra
  \item May have used AI tools like ChatGPT or GitHub Copilot
  \item Does \emph{not} necessarily know NLP, linguistics, CUDA, or distributed systems
\end{itemize}

If some of those prerequisites feel shaky, Part~0 provides self-contained refreshers on everything you need.

\section*{How to Read This Book}

The book is organized into six parts plus appendices:
\begin{description}
  \item[Part 0: Prerequisites] Probability, linear algebra, PyTorch, GPU basics, and linguistics background
  \item[Part I: Foundations] The language modeling problem, tokenization, and the path from n-grams to neural models
  \item[Part II: Transformer Internals] Attention, Transformer blocks, normalization, and architectural variants
  \item[Part III: Systems] Resource accounting, GPUs, kernels, parallelism, and inference
  \item[Part IV: Scaling and Evaluation] Scaling laws, model recipes, and benchmarking
  \item[Part V: Data] Data sources, filtering, deduplication, and mixing
  \item[Part VI: Alignment] Instruction tuning, RLHF, DPO, reasoning RL, and safety
\end{description}

You can read sequentially, or follow one of the suggested reading paths in Chapter~0.

\section*{Neuroscience Analogies}

Throughout this book, we use analogies from neuroscience and brain science to build intuition. Every analogy box clearly states where the analogy helps, where it breaks down, and what intuition you should keep. We never imply that language models are brains---but the parallels can sharpen your understanding of both.

\section*{Acknowledgments}

This book builds on the materials, lectures, and assignments created by Tatsunori Hashimoto, Percy Liang, and the CS336 teaching staff. We are grateful for their vision of teaching language modeling from scratch.
